% Part: normal-modal-logic
% Chapter: filtrations
% Section: S5-decidable

\documentclass[../../../include/open-logic-section]{subfiles}

\begin{document}

\olfileid[tr]{nml}{fil}{dec}

\olsection{\Log{S5} Karar Verilebilirdir}
Sonlu model özelliği, şemalarla verilen modal mantık sistemlerinin
\emph{karar verilebilir} olduğunu, yani !!a{formula}{}ün sistem içinde
!!{derivable} olup olmadığını belirleyen hesaplanabilir bir yordam
bulunduğunu göstermenin kolay bir yolunu verir.

\begin{thm}
  \Log{S5} karar verilebilirdir.
\end{thm}

\begin{proof}
  $!A$ verilsin ve $!A$ içinde geçen !!{propositional variable}s{}nin
  $p_1$,\dots,$p_k$ arasında bulunduğunu varsayalım. Her $n$ için
  $p_1$,\dots,$p_k$'ye değer atayan $n$ dünyalı yalnızca sonlu sayıda model
  olduğundan, \Log{S5}'in bütün teoremlerini sistematik bir biçimde
  ispatlar üreterek ve $1$, $2$, \dots dünyalı bütün evrensel modelleri üretip
  $!A$'nın bunlardan birinin bir dünyasında yanlış olup olmadığını denetleyerek
  \emph{paralel olarak} numaralandırabiliriz. İki paralel süreçten biri
  sonunda yanıt verecektir; çünkü
  \olref[com][fra]{thm:generaldet} ile \olref[fmp]{cor:S5fmp} gereğince ya
  $!A$ !!{derivable}{}dir ya da sonlu bir evrensel modelde yanlıştır.
\end{proof}

Yukarıdaki ispat \Log{S5} için işler; çünkü evrensel modellerin süzmeleri
kendiliğinden evrenseldir. Aynısı yansımalılık ve serilik için de geçerlidir;
öteki özellikler için daha çok çalışma gerekir.

\end{document}
